\clearpage
\section*{ВВЕДЕНИЕ}
\addcontentsline{toc}{section}{ВВЕДЕНИЕ}

Современные финансовые рынки характеризуются высокой скоростью обновления информации, существенной изменчивостью цен и постоянной конкуренцией между участниками, использующими различные аналитические и вычислительные методы. В этой среде задача прогнозирования биржевых котировок остается одной из ключевых как для прикладной финансовой аналитики, так и для области искусственного интеллекта. Однако накопленный опыт показывает, что сама по себе задача прогноза цены не исчерпывает практическую постановку проблемы. Даже при наличии модели, демонстрирующей приемлемое качество по статистическим критериям, ее полезность для принятия торговых решений может оказаться ограниченной. Модель может корректно предсказывать направление части движений, но при этом формировать чрезмерно частые сигналы, создавать высокую просадку или не учитывать транзакционные издержки [3, 6, 7].

Актуальность темы определяется сразу несколькими обстоятельствами. Во-первых, финансовые временные ряды являются типичным примером сложных нестационарных данных, для которых недостаточно только классических статистических методов. Во-вторых, широкое распространение получили модели машинного обучения и глубокого обучения, ориентированные на выделение сложных зависимостей в данных. В-третьих, в последние годы заметно вырос интерес к обучению с подкреплением как к подходу, который позволяет не просто прогнозировать цену, а обучать агента непосредственному выбору действий на рынке. Для задач алгоритмической торговли это особенно важно, поскольку конечная цель практической системы состоит не в минимизации ошибки прогноза как таковой, а в повышении эффективности последовательных торговых решений с учетом риска, просадки, комиссий и устойчивости поведения на различных рыночных режимах [3, 7, 11, 12, 21].

Дополнительную значимость теме придает замечание, сделанное научным руководителем: при обосновании актуальности необходимо не только описать применимость RL-подходов, но и провести сравнение существующих аналогов, а также четко показать проблему, которую решает работа. Действительно, без такого сравнения исследование рискует превратиться в изолированное описание одного алгоритма. Поэтому в настоящей работе обучение с подкреплением рассматривается не как самоцель, а как один из классов методов, сопоставляемый с традиционными baseline-стратегиями, статистическими моделями и классическими алгоритмами машинного обучения в рамках единого экспериментального стенда.

Проблема исследования состоит в следующем. Большинство подходов к анализу биржевых данных оцениваются либо по ошибке прогноза, либо по отдельным финансовым показателям, но не всегда в рамках единой методологии. На практике это приводит к нескольким методологическим трудностям. Во-первых, сравнение моделей оказывается некорректным, если они обучаются и тестируются на разных интервалах или по разным правилам формирования сигналов. Во-вторых, высокая классификационная точность не эквивалентна высокой доходности стратегии, поскольку итоговый результат зависит от частоты сделок, величины прибыли на одну сделку, транзакционных издержек и характера рыночного режима. В-третьих, многие работы рассматривают RL-агентов обособленно, без сопоставления с простыми и интерпретируемыми альтернативами, такими как Buy \& Hold или стратегия на пересечении скользящих средних [6, 7, 27, 28].

Таким образом, возникает необходимость в разработке программного решения, которое позволяет на едином наборе биржевых данных формировать признаки, обучать модели, запускать алгоритмы принятия торговых решений, выполнять backtesting и рассчитывать показатели качества по единым правилам. Именно создание и экспериментальная проверка такого решения составляют центральное содержание настоящей выпускной квалификационной работы.

Целью работы является разработка, программная реализация и экспериментальная проверка программного решения для прогнозирования и оптимизации торговых решений на основе биржевых котировок с использованием baseline-стратегий, моделей машинного обучения и методов обучения с подкреплением.

Для достижения поставленной цели в работе решаются следующие задачи:
\begin{enumerate}[leftmargin=1.25cm,label=\arabic*)]
\item проанализировать особенности биржевых котировок как финансовых временных рядов и определить требования к разрабатываемому программному решению;
\item выполнить обзор существующих подходов и аналогов, необходимых для выбора состава реализуемых модулей и методов сравнения;
\item спроектировать архитектуру программного решения, включающего загрузку данных, формирование признаков, модуль принятия решений и подсистему оценки результатов;
\item реализовать модуль подготовки OHLCV-данных, модуль генерации признаков, baseline-стратегии, классические ML-модели и RL-среду для обучения агентов DQN и PPO;
\item формализовать задачу обучения с подкреплением, включая описание состояния среды, пространства действий, функции вознаграждения и ограничений торговой логики;
\item реализовать процедуры обучения, тестирования и backtesting, обеспечивающие воспроизводимый запуск вычислительного эксперимента;
\item провести экспериментальную проверку разработанного программного решения и сравнить реализованные подходы по классификационным и финансовым метрикам;
\item выполнить анализ полученных результатов, определить ограничения решения и наметить направления его дальнейшего развития.
\end{enumerate}

Объектом исследования являются методы анализа и принятия решений на основе биржевых котировок.

Предметом исследования выступают алгоритмы прогнозирования и оптимизации торговых стратегий на финансовых временных рядах, в том числе baseline-подходы, модели машинного обучения и методы обучения с подкреплением.

В работе применяются методы статистического анализа временных рядов, методы машинного обучения, методы обучения с подкреплением, методы финансовой оценки торговых стратегий, а также методы программной инженерии для построения воспроизводимого вычислительного контура. Практическая реализация выполнена на языке Python с использованием библиотек \textit{pandas}, \textit{NumPy}, \textit{scikit-learn}, \textit{yfinance}, \textit{Gymnasium}, \textit{stable-baselines3} и \textit{matplotlib} [15, 16, 17, 19, 20, 24].

Практическая значимость работы заключается в создании работоспособного программного решения, пригодного для воспроизводимого запуска экспериментов по алгоритмической торговле на исторических котировках. Разработанный pipeline позволяет автоматически загружать данные, формировать признаки, обучать реализованные модели, запускать стратегии на тестовом интервале, рассчитывать финансовые показатели и сохранять результаты в виде отчетных артефактов. Такой подход представляет интерес как для учебных целей, так и для дальнейшего наращивания функциональности за счет новых признаков, алгоритмов и сценариев оценки.

Научная новизна работы в рамках учебного исследования состоит в том, что RL-подход рассматривается не изолированно, а как часть разработанного программного контура, в котором прогнозирующие модели и торговые стратегии сравниваются по единым правилам входных данных, тестирования и расчета метрик. Это позволяет сместить акцент с абстрактного качества прогноза на более прикладной вопрос: может ли реализованный RL-агент улучшить не только точность сигнала, но и риск-скорректированную результативность торговой системы.

Структурно работа включает введение, четыре главы, заключение, список использованных источников и приложение. В первой главе рассматриваются предметная область, существующие подходы и аналоги, а также формулируется проблема исследования. Во второй главе описываются проектирование метода, подготовка данных, признаки, постановка RL-задачи и система метрик. Третья глава посвящена реализации программного решения и экспериментального стенда. В четвертой главе представлены результаты экспериментального моделирования, их сопоставление и интерпретация. В заключении сформулированы основные результаты, ограничения и направления дальнейшего развития работы.
